\documentclass{article}
\usepackage{geometry}
\geometry{margin=1.5cm, vmargin={0pt,1cm}}
\setlength{\topmargin}{-1cm}
\setlength{\headheight}{12.5pt}
\setlength{\paperheight}{29.7cm}
\setlength{\textheight}{25.3cm}

% useful packages.
\usepackage[UTF8]{ctex}
\usepackage{fancyhdr}
\usepackage{layout}
\usepackage{luacode}
\usepackage{tikz}

% import common macros
% % % % % % % % % % % % % % % % % % % % % % % % % % % % % % % %
% Common Macros for LaTeX
% % % % % % % % % % % % % % % % % % % % % % % % % % % % % % % %

% Useful packages
\usepackage{amsfonts}
\usepackage{amsmath}
\usepackage{amssymb}
\usepackage{amsthm}
\usepackage{enumerate}
\usepackage{graphicx}
\usepackage{multicol}
\usepackage[ruled,linesnumbered]{algorithm2e}

% Math operators and common symbols
\newcommand{\dif}{\mathrm{d}}
\newcommand{\innerProd}[1]{\left\langle #1 \right\rangle}
\newcommand{\difFrac}[2]{\frac{\dif #1}{\dif #2}}
\newcommand{\pdfFrac}[2]{\frac{\partial #1}{\partial #2}}
\newcommand{\T}{\mathrm{T}}
\newcommand{\norm}[1]{\left\| #1 \right\|}
\newcommand{\abs}[1]{\left| #1 \right|}

% Vector and Matrix shortcuts
\newcommand{\vecTwo}[2]{
  \begin{bmatrix}
    #1 \\
    #2
\end{bmatrix}}
\newcommand{\vecThree}[3]{
  \begin{bmatrix}
    #1 \\
    #2 \\
    #3
\end{bmatrix}}
\newcommand{\vecTwoH}[2]{
  \begin{bmatrix}
    #1 & #2
  \end{bmatrix}
}
\newcommand{\vecThreeH}[3]{
  \begin{bmatrix}
    #1 & #2 & #3
  \end{bmatrix}
}

% Common fields
\newcommand{\R}{\mathbb{R}}
\newcommand{\C}{\mathbb{C}}
\newcommand{\Q}{\mathbb{Q}}
\newcommand{\Z}{\mathbb{Z}}
\newcommand{\N}{\mathbb{N}}

% Common Operators
\renewcommand{\Re}{\mathrm{Re}}
\renewcommand{\Im}{\mathrm{Im}}
\newcommand{\Arg}{\mathrm{Arg}}
\newcommand{\Log}{\mathrm{Log}}

\newcommand{\arccot}{\mathrm{arccot}}
\newcommand{\arcsec}{\mathrm{arcsec}}
\newcommand{\arccsc}{\mathrm{arccsc}}

% Theorem-like environments
\theoremstyle{definition}
\newtheorem{theorem}{Theorem}[section]
\newtheorem{lemma}[theorem]{Lemma}
\newtheorem{corollary}[theorem]{Corollary}
\newtheorem{definition}[theorem]{Definition}
\newtheorem{remark}[theorem]{Remark}
\newtheorem{exercise}{Exercise}[section]

% Support for individual Chinese characters using fontspec (for LuaLaTeX)
\usepackage{fontspec}
\newfontfamily\zhfont{SimSun} % Search system fonts by name
\newcommand{\zhstr}[1]{{\zhfont #1}}

% Solutions and proofs
\AtBeginDocument{
  \ifdefined\ctexset
  \newenvironment{solution}{
  \begin{proof}[解]}{
  \end{proof}}
  \else
  \newenvironment{solution}{
  \begin{proof}[Solution]}{
  \end{proof}}
  \fi
}

\usepackage{luacode}
% Utility to get environment variables (requires lualatex)
\newcommand{\getEnv}[2]{\directlua{tex.print(os.getenv("#1") or "#2")}}


\usepackage{listings}
\usepackage{xcolor}

% C++ 源码的挂载配置
\lstset{
  language=C++,
  basicstyle=\ttfamily\small,
  keywordstyle=\bfseries\color{blue},
  commentstyle=\color{gray},
  stringstyle=\color{red},
  breaklines=true,
  showstringspaces=false,
  tabsize=4,
  frame=single,
  framesep=5pt,
  rulecolor=\color{gray},
  captionpos=b,
  numberstyle=\tiny\color{gray},
  numbers=left,
}

\newcommand{\fl}{\operatorname{fl}}
\newcommand{\UFL}{\operatorname{UFL}}
\newcommand{\OFL}{\operatorname{OFL}}
\newcommand{\F}{\mathbb{F}}

% Name and uID
\newcommand{\name}{\getEnv{NAME}{NAME}}
\newcommand{\uid}{\getEnv{UID}{UID}}
\newcommand{\course}{数值分析}
\newcommand{\hwNum}{1}

\begin{document}

\pagestyle{fancy}
\fancyhead{}
\lhead{\name (\uid)}
\chead{\course \#\hwNum}
\rhead{\today}

\section{习题 1.6}

\subsection*{I}

把十进制整数 477 化为 $\beta = 2$ 的规格化 FPN。

\begin{solution}
  除 2 取余得
  \[
    477 = (111011101)_2 = 1.11011101 \times 2^{8}.
  \]
\end{solution}

\subsection*{II}

把十进制分数 $3/5$ 化为 $\beta = 2$ 的规格化 FPN。

\begin{solution}
  乘 2 取整，小数部分以 4 为周期：
  \[
    \frac{3}{5} = (0.\overline{1001})_2 = 1.\overline{1001} \times 2^{-1}.
  \]
\end{solution}

\subsection*{III}

设 $x = \beta^{e}$ 为规格化 FPN，其中 $e \in \Z$ 且 $L < e < U$，$x_L, x_R \in \F$
为与 $x$ 相邻的两个规格化 FPN 且 $x_L < x < x_R$。证明
$x_R - x = \beta(x - x_L)$。

\begin{proof}
  $x = \beta^{e}$ 是 binade $[\beta^{e}, \beta^{e+1})$ 的左端点，故 $x_L$ 是前一个
  binade 中尾数最大的元素，$x_R$ 是本文 binade 中尾数为 $1 + \beta^{1-p}$ 的元素：
  \[
    x_L = (\beta - \beta^{1-p})\beta^{e-1}, \qquad
    x_R = (1 + \beta^{1-p})\beta^{e}.
  \]
  于是
  \[
    x - x_L = \beta^{e} - (\beta - \beta^{1-p})\beta^{e-1} = \beta^{e-p},
  \]
  \[
    x_R - x = (1 + \beta^{1-p})\beta^{e} - \beta^{e} = \beta^{e+1-p} = \beta(x - x_L).
  \]
\end{proof}

\subsection*{IV}

利用 II 的结果求 IEEE 754 单精度下 $x = 3/5$ 的两个相邻规格化 FPN，并给出 $\fl(x)$
与相对舍入误差。

\begin{solution}
  单精度 $p = 24$，尾数取 23 位小数。由 II，$x$ 的第 24 位尾数及其后不全为零，
  两个相邻的 FPN 为
  \[
    x_L = (1.\underbrace{00110011001100110011001}_{23})_2 \times 2^{-1}
        \approx 0.5999999642372131,
  \]
  \[
    x_R = (1.\underbrace{00110011001100110011010}_{23})_2 \times 2^{-1}
        \approx 0.6000000238418579.
  \]
  二者相距一个 ulp，即 $x_R - x_L = 2^{-24}$。第 24 位尾数之后余下的部分为
  $(0.\overline{1001})_2 = 3/5$，故
  \[
    x - x_L = \frac{3}{5} \cdot 2^{-24}, \qquad
    x_R - x = \frac{2}{5} \cdot 2^{-24},
  \]
  $x$ 距 $x_R$ 更近，所以
  \[
    \fl(x) = x_R = 1.00110011001100110011010 \times 2^{-1}.
  \]
  相对舍入误差为
  \[
    \left| \frac{\fl(x) - x}{x} \right|
    = \frac{(2/5)2^{-24}}{3/5}
    = \frac{2}{3} \cdot 2^{-24}
    \approx 3.97 \times 10^{-8}
    = \frac{2}{3}\epsilon_u.
  \]
\end{solution}

\subsection*{V}

若 IEEE 754 单精度协议不采用就近舍入，而是直接截断多余位，单位舍入误差是多少？

\begin{solution}
  截断时误差不超过一个 ulp，故
  \[
    |\fl(x) - x| < \epsilon_M |x| = \beta^{1-p} |x|.
  \]
  而 Definition 1.28 中 $\epsilon_u = \tfrac{1}{2}\epsilon_M$，所以截断下的单位舍入
  误差为
  \[
    \epsilon_u = \epsilon_M = 2^{1-24} = 2^{-23},
  \]
  是就近舍入的两倍。
\end{solution}

\subsection*{VI}

对 Example 1.34 重做，取 $b = 8.769 \times 10^{4}$、$b = -5.678 \times 10^{0}$、
$b = -5.678 \times 10^{3}$。

\begin{solution}
  $F: (10, 4, -7, 8)$，$a = 1.234 \times 10^{4}$，$\epsilon_u(p) = 0.5 \times 10^{-3}$。

  \textbf{$b = 8.769 \times 10^{4}$。} $|b| > |a|$，先交换两个操作数使 $|a| \ge |b|$。
  此时 $e_a - e_b = 0$，对齐后 $M_c = 8.769 + 1.234 = 10.003$，$e_c = 4$。
  因 $|M_c| \in (10, 100)$，执行 (iii) 中的正规化：$M_c \leftarrow 1.0003$，
  $e_c \leftarrow 5$。舍入到 4 位有效数字，第 5 位为 3 故舍去，得
  \[
    \fl(a + b) = 1.000 \times 10^{5}.
  \]

  \textbf{$b = -5.678 \times 10^{0}$。} $e_a - e_b = 4 \le p + 1$，对齐得
  $M_b = -0.0005678$，于是
  \[
    M_c = 1.234 - 0.0005678 = 1.2334322 \in [1, 10),
  \]
  无需正规化。舍入到 4 位有效数字，第 5 位为 4 故舍去，得
  \[
    \fl(a + b) = 1.233 \times 10^{4}.
  \]

  \textbf{$b = -5.678 \times 10^{3}$。} 对齐得 $M_b = -0.5678$，于是
  \[
    M_c = 1.234 - 0.5678 = 0.6662 \in (0, \beta - \epsilon_u(p)),
  \]
  触发 (iii) 中的左移正规化：$M_c \leftarrow 6.662$，$e_c \leftarrow 3$。
  舍入到 4 位不变，得
  \[
    \fl(a + b) = 6.662 \times 10^{3}.
  \]
\end{solution}

\subsection*{VII}

设两个 FPN 的除法在精度为 $2p$ 的寄存器中完成，举一例说明机器算术模型的结论不再成立。

\begin{solution}
  取 $F(2, 3, -1, 1)$ 中的
  \[
    a = 1.00 \times 2^{0} = 1, \qquad b = 1.11 \times 2^{0} = \frac{7}{4},
    \qquad \epsilon_u = \frac{1}{8}.
  \]
  真值 $a/b = 4/7 = (0.100100100\ldots)_2$。商小于 1，故按 Definition 1.40 的顺序，
  寄存器先按第 (ii) 步做除法，再按第 (iii) 步正规化。商的小数点前那个 0 占去一个
  数位，$2p = 6$ 位的寄存器只剩 $2^{-5}$ 可用：
  \[
    \frac{4}{7} = (0.100100100\ldots)_2 \longrightarrow (0.10010)_2 = \frac{9}{16},
  \]
  余下的位不足半个 ulp，故向下舍入。左移正规化得
  $M_c = 2 \times \frac{9}{16} = \frac{9}{8} = (1.001)_2$，$e_c = -1$。
  第 (v) 步舍入到 $p = 3$ 位：$\frac{9}{8}$ 与 $1.00$、$1.01$ 的距离都是 $0.125$，
  恰好平局，按 round to even 取 $1.00$，于是
  \[
    \fl\left(\frac{a}{b}\right) = 1.00 \times 2^{-1} = \frac{1}{2},
    \qquad
    |\delta| = \left| \frac{1/2}{4/7} - 1 \right| = \frac{1}{8} = \epsilon_u.
  \]
  Theorem 1.42 要求 $|\delta| < \epsilon_u$，此处取到等号，模型不成立。

  若寄存器取 $2p + 1 = 7$ 位，$4/7$ 舍入到最近的格点得
  $(0.100101)_2 = \frac{37}{64}$，正规化后 $M_c = 1.15625$，舍入到 3 位得
  $1.01 \times 2^{-1} = 5/8$，$|\delta| = \frac{3}{32} < \epsilon_u$，模型重新成立。
\end{solution}

\subsection*{VIII}

把正数 $a_i > 0$ 按大小升序相加可以减小 Theorem 1.43 中的误差项 $\delta$，举两例。

\begin{solution}
  当 $a_{k+1} \ll s_k$ 时，对齐后 $a_{k+1}$ 的各位落在大数的 ulp 之下而被整个舍掉。
  升序时小数先聚成较大的部分和，再加到大数上就不会被吞掉。以下两例都在
  $F(10, 4, -7, 8)$ 中计算，每步部分和舍入到 4 位有效数字。

  \textbf{例 1.} 计算 $10^{4} + \underbrace{1 + \cdots + 1}_{10}$。

  降序：每次加 1 都因对齐被舍去，最终得 $1.000 \times 10^{4}$，误差 10。

  升序：十个 1 先累加得 10，再算 $10^{4} + 10 = 10010$，保留 4 位得
  $1.001 \times 10^{4}$，即精确值。

  \textbf{例 2.} 计算 $10^{4} + \underbrace{6 + \cdots + 6}_{8}$，精确值为 10048。

  降序：每步 $+6$ 都舍入到最近的 10 的倍数，最终得 $1.008 \times 10^{4}$，
  误差 $+32$。

  升序：八个 6 先累加得 48，再算 $10048$，舍入得 $1.005 \times 10^{4}$，误差 $+2$。
\end{solution}

\subsection*{IX}

设 $a_i, b_i \in \F$，推导 $\fl(a_1 b_1 + a_2 b_2 + a_3 b_3)$，并讨论
$\fl(\sum_i \prod_j a_{i,j})$ 相应的推导。

\begin{solution}
  记 $p_i := \fl(a_i b_i)$。由 Lemma 1.39，乘积在 $2p$ 位寄存器中精确，舍入只在末尾：
  \[
    p_i = a_i b_i (1 + \delta_i), \qquad |\delta_i| < \epsilon_u.
  \]
  两次加法各引入一个误差因子，由 Lemma 1.36
  \[
    \fl(p_1 + p_2) = (p_1 + p_2)(1 + \eta_1), \qquad
    \fl\big(\fl(p_1 + p_2) + p_3\big) = \big(\fl(p_1 + p_2) + p_3\big)(1 + \eta_2),
  \]
  $|\eta_1|, |\eta_2| < \epsilon_u$。代入 $p_i$ 得
  \[
    \fl(a_1b_1 + a_2b_2 + a_3b_3)
    = \sum_{i=1}^{3} a_i b_i (1 + \theta_i),
    \qquad |\theta_i| \le 3\epsilon_u + 9\epsilon_u^{2} = 3\epsilon_u + O(\epsilon_u^2),
  \]
  其中前两项含 3 个因子、末项含 2 个。每个项最多经历 3 次舍入：乘积一次，两次加法
  各一次。误差阶为 $3\epsilon_u$。

  对一般的 $\fl(\sum_i \prod_j a_{i,j})$，设外层求和有 $m$ 项、每个内积含 $k$ 个因子，
  则内积要 $k - 1$ 次舍入、外层对每项要至多 $m - 1$ 次，误差最大的项含 $k + m - 2$ 个
  $(1 + \epsilon)$ 因子，
  \[
    |\theta_i| \lesssim (k + m - 2)\epsilon_u,
  \]
  舍入误差随运算总数线性累积，与 Theorem 1.43 的
  $|\delta_n| < (1 + \epsilon_u)^n - 1 \approx n\epsilon_u$ 一致。
\end{solution}

\subsection*{X}

计算 $x = \frac{1}{4}$ 时减法 $1 - \cos x$ 损失多少位精度。

\begin{solution}
  $\cos(1/4) \approx 0.9689124217$，故
  \[
    1 - \cos\frac{1}{4} \approx 0.0310875783 \in [2^{-6}, 2^{-5}).
  \]
  对 Theorem 1.46 的条件 (1.24) 取 $x = 1$、$y = \cos(1/4)$、$\beta = 2$，由
  \[
    2^{-6} \le 1 - \frac{y}{x} \le 2^{-5}
  \]
  得 $s = 5$、$t = 6$，即损失的最高有效位至少 5 位、至多 6 位。
\end{solution}

\subsection*{XI}

给出至少两种计算 $1 - \cos x$ 以避免灾难性抵消的方法。

\begin{solution}
  \textbf{方法一。} 用半角公式
  \[
    1 - \cos x = 2 \sin^{2} \frac{x}{2},
  \]
  当 $x$ 小时 $\sin(x/2)$ 与 $x/2$ 同量级，不再有相近数相减。

  \textbf{方法二。} 分子有理化
  \[
    1 - \cos x = \frac{\sin^{2} x}{1 + \cos x},
  \]
  分母 $1 + \cos x$ 接近 2，而分子是乘积不是减法。

  \textbf{方法三。} 对小的 $x$ 用 Taylor 展开
  \[
    1 - \cos x = \frac{x^{2}}{2!} - \frac{x^{4}}{4!} + \frac{x^{6}}{6!} - \cdots
  \]
\end{solution}

\section{习题 1.7}

考虑特征为 $\beta = 2$、$p = 3$、$L = -1$、$U = +1$ 的规格化 FPN 系统 $F$。

\begin{solution}
  \textbf{UFL 与 OFL。} 由 Definition 1.19，
  $\UFL(F) = \beta^{L} = 2^{-1} = 0.5$，
  $\OFL(F) = \beta^{U}(\beta - \beta^{1-p}) = 2(2 - 2^{-2}) = 3.5$，
  单位舍入误差 $\epsilon_u = 2^{-2}/2 = 0.125$。

  \textbf{$F$ 的全部元素。} 枚举得正的规格化元素如下：
  \begin{center}
    % Generated by hw1/main.cpp -- do not edit by hand.
\begin{tabular}{cc}
\hline
值 & 规格化形式 \\
\hline
0.5 & $1.00 \times 2^{-1}$ \\
0.625 & $1.01 \times 2^{-1}$ \\
0.75 & $1.10 \times 2^{-1}$ \\
0.875 & $1.11 \times 2^{-1}$ \\
1 & $1.00 \times 2^{0}$ \\
1.25 & $1.01 \times 2^{0}$ \\
1.5 & $1.10 \times 2^{0}$ \\
1.75 & $1.11 \times 2^{0}$ \\
2 & $1.00 \times 2^{1}$ \\
2.5 & $1.01 \times 2^{1}$ \\
3 & $1.10 \times 2^{1}$ \\
3.5 & $1.11 \times 2^{1}$ \\
\hline
\end{tabular}

  \end{center}
  共 12 个；计入相反数与 0，$F$ 共 25 个元素。Corollary 1.21 给出
  \[
    \#F = 2^{p}(U - L + 1) + 1 = 2^{3}\bigl(1 - (-1) + 1\bigr) + 1 = 25,
  \]
  与枚举一致。

  $F$ 在实轴上的位置：
  \begin{center}
    % Generated by hw1/main.cpp -- do not edit by hand.
\begin{tikzpicture}[xscale=1.9, yscale=0.8]
  \draw[thick,<->] (-3.9,0) -- (3.9,0);
  % elements of F
  \draw[red] (0.5,0.18) -- (0.5,-0.18);
  \draw[red] (-0.5,0.18) -- (-0.5,-0.18);
  \draw[red] (0.625,0.18) -- (0.625,-0.18);
  \draw[red] (-0.625,0.18) -- (-0.625,-0.18);
  \draw[red] (0.75,0.18) -- (0.75,-0.18);
  \draw[red] (-0.75,0.18) -- (-0.75,-0.18);
  \draw[red] (0.875,0.18) -- (0.875,-0.18);
  \draw[red] (-0.875,0.18) -- (-0.875,-0.18);
  \draw[red] (1,0.18) -- (1,-0.18);
  \draw[red] (-1,0.18) -- (-1,-0.18);
  \draw[red] (1.25,0.18) -- (1.25,-0.18);
  \draw[red] (-1.25,0.18) -- (-1.25,-0.18);
  \draw[red] (1.5,0.18) -- (1.5,-0.18);
  \draw[red] (-1.5,0.18) -- (-1.5,-0.18);
  \draw[red] (1.75,0.18) -- (1.75,-0.18);
  \draw[red] (-1.75,0.18) -- (-1.75,-0.18);
  \draw[red] (2,0.18) -- (2,-0.18);
  \draw[red] (-2,0.18) -- (-2,-0.18);
  \draw[red] (2.5,0.18) -- (2.5,-0.18);
  \draw[red] (-2.5,0.18) -- (-2.5,-0.18);
  \draw[red] (3,0.18) -- (3,-0.18);
  \draw[red] (-3,0.18) -- (-3,-0.18);
  \draw[red] (3.5,0.18) -- (3.5,-0.18);
  \draw[red] (-3.5,0.18) -- (-3.5,-0.18);
  % axis labels
  \node[below] at (-3,-0.25) {$-3$};
  \node[below] at (-2,-0.25) {$-2$};
  \node[below] at (-1,-0.25) {$-1$};
  \node[below] at (0,-0.25) {$0$};
  \node[below] at (1,-0.25) {$1$};
  \node[below] at (2,-0.25) {$2$};
  \node[below] at (3,-0.25) {$3$};
\end{tikzpicture}

  \end{center}

  \textbf{次正规数。} 由 Definition 1.15，次正规数是 $e = L$ 且尾数 $m \in (0,1)$
  的 FPN，共 $\beta^{p-1} - 1 = 3$ 个：
  \begin{center}
    % Generated by hw1/main.cpp -- do not edit by hand.
\begin{tabular}{cc}
\hline
值 & 形式 \\
\hline
0.125 & $0.01 \times 2^{-1}$ \\
0.25 & $0.10 \times 2^{-1}$ \\
0.375 & $0.11 \times 2^{-1}$ \\
\hline
\end{tabular}

  \end{center}

  \textbf{扩充后的 $F$。}
  \begin{center}
    % Generated by hw1/main.cpp -- do not edit by hand.
\begin{tikzpicture}[xscale=1.9, yscale=0.8]
  \draw[thick,<->] (-3.9,0) -- (3.9,0);
  % elements of F
  \draw[red] (0.5,0.18) -- (0.5,-0.18);
  \draw[red] (-0.5,0.18) -- (-0.5,-0.18);
  \draw[red] (0.625,0.18) -- (0.625,-0.18);
  \draw[red] (-0.625,0.18) -- (-0.625,-0.18);
  \draw[red] (0.75,0.18) -- (0.75,-0.18);
  \draw[red] (-0.75,0.18) -- (-0.75,-0.18);
  \draw[red] (0.875,0.18) -- (0.875,-0.18);
  \draw[red] (-0.875,0.18) -- (-0.875,-0.18);
  \draw[red] (1,0.18) -- (1,-0.18);
  \draw[red] (-1,0.18) -- (-1,-0.18);
  \draw[red] (1.25,0.18) -- (1.25,-0.18);
  \draw[red] (-1.25,0.18) -- (-1.25,-0.18);
  \draw[red] (1.5,0.18) -- (1.5,-0.18);
  \draw[red] (-1.5,0.18) -- (-1.5,-0.18);
  \draw[red] (1.75,0.18) -- (1.75,-0.18);
  \draw[red] (-1.75,0.18) -- (-1.75,-0.18);
  \draw[red] (2,0.18) -- (2,-0.18);
  \draw[red] (-2,0.18) -- (-2,-0.18);
  \draw[red] (2.5,0.18) -- (2.5,-0.18);
  \draw[red] (-2.5,0.18) -- (-2.5,-0.18);
  \draw[red] (3,0.18) -- (3,-0.18);
  \draw[red] (-3,0.18) -- (-3,-0.18);
  \draw[red] (3.5,0.18) -- (3.5,-0.18);
  \draw[red] (-3.5,0.18) -- (-3.5,-0.18);
  % subnormal numbers (Definition 1.15)
  \draw[blue] (0.125,0.18) -- (0.125,-0.18);
  \draw[blue] (-0.125,0.18) -- (-0.125,-0.18);
  \draw[blue] (0.25,0.18) -- (0.25,-0.18);
  \draw[blue] (-0.25,0.18) -- (-0.25,-0.18);
  \draw[blue] (0.375,0.18) -- (0.375,-0.18);
  \draw[blue] (-0.375,0.18) -- (-0.375,-0.18);
  % axis labels
  \node[below] at (-3,-0.25) {$-3$};
  \node[below] at (-2,-0.25) {$-2$};
  \node[below] at (-1,-0.25) {$-1$};
  \node[below] at (0,-0.25) {$0$};
  \node[below] at (1,-0.25) {$1$};
  \node[below] at (2,-0.25) {$2$};
  \node[below] at (3,-0.25) {$3$};
\end{tikzpicture}

  \end{center}
  红色刻度是 $F$ 的规格化元素，蓝色刻度是次正规数。它们与 $\UFL = 0.5$ 的间距都是
  最小间距 $2^{-3}$，即 Definition 1.27 所说的 gradual underflow。
  扩充后 $F$ 共有 $25 + 2 \times 3 = 31$ 个元素。
\end{solution}

程序的执行结果即上文各表格与插图。核心源码如下：\texttt{fpn.hpp} 与
\texttt{fpn.cpp} 是 FPN 系统的实现，\texttt{main.cpp} 是本题的程序。
表格打印等辅助函数位于 \texttt{common/} 下，与本题算法无关，不再列出。

\lstinputlisting[language=C++, caption={作业程序 main.cpp}]{main.cpp}
\lstinputlisting[language=C++, caption={FPN 系统 fpn.hpp}]{../common/include/na/fpn.hpp}
\lstinputlisting[language=C++, caption={FPN 系统 fpn.cpp}]{../common/src/fpn.cpp}

\end{document}
